\paragraph{LLM 路由。}
具有不同成本--能力特征的 LLM 大量涌现，推动了智能模型选择的研究。FrugalGPT~\citep{chen2023frugalgpt}将模型从廉价到昂贵依次级联，直至达到置信度阈值。EmbedLLM~\citep{wang2024embedllm}学习嵌入，以预测模型在特定查询上的表现。RouterDC~\citep{chen2024routerdc}使用双重对比学习训练有效的路由器。Avengers~\cite{zhang2025avengers}使用简单的聚类式路由方案编排十个 7B 模型池，实现超过 GPT-4.1 的性能；AvengersPro~\cite{zhang2025beyond}则将这一设计扩展为在平衡评测设置下提供帕累托最优的性能--成本权衡。这些方法聚焦于\emph{单轮}路由：给定一个查询，选择一个模型作答。面向多轮场景，Router-R1~\citep{zhang2025router}训练策略 LLM，通过强化学习交错进行推理和路由。不同于依赖重型 LLM 路由器的 Router-R1，我们的方法在联合历史--模型嵌入上学习\textit{轻量级}结果估计器，实现高成本效益的轮级路由。

\paragraph{LLM 工具使用。}
由工具增强的 LLM 能够完成复杂的多步任务~\citep{schick2023toolformer,qin2023toolllm}。ReAct~\citep{yao2023react}在单个提示词内交错推理与动作；Toolformer~\citep{schick2023toolformer}则微调模型以自主调用工具。近期研究已在多种领域探索基于 LLM 的工具使用，包括软件工程~\citep{jimenez2024swebench,yang2024sweagent}、网页浏览~\citep{zhou2024webarena}、操作系统交互~\citep{xie2024osworld}、复杂推理~\citep{mialon2023gaiabenchmarkgeneralai}与自主研究~\citep{team2025tongyi}。这些任务本质上具有交互性，通常无法在单轮内解决，因而需要重复的动作--观察循环。多数工具使用框架依赖固定的底层模型；我们的工作引入一个在模型间动态切换的路由层，将模型选择视为每一步的自适应决策。这一方法与现有工具使用方法互补，并可推广至多轮 LLM 系统。
